\documentclass[a4paper,11pt]{article}

% Packages for encoding and formatting
\usepackage[utf8]{inputenc}
\usepackage[T1]{fontenc}
\usepackage{geometry}
\usepackage{xcolor}
\usepackage{hyperref}
\usepackage{listings}
\usepackage{enumitem}
\usepackage{titlesec}

% Page margins
\geometry{margin=1in}

% Hyperlink setup
\hypersetup{
    colorlinks=true,
    linkcolor=blue,
    filecolor=magenta,      
    urlcolor=blue,
}

% Code listing style (Terminal-like)
\definecolor{codegray}{rgb}{0.95,0.95,0.95}
\definecolor{codeblack}{rgb}{0,0,0}

\lstdefinestyle{mystyle}{
    backgroundcolor=\color{codegray},   
    basicstyle=\ttfamily\small\color{codeblack},
    breakatwhitespace=false,         
    breaklines=true,                 
    captionpos=b,                    
    keepspaces=true,                 
    numbers=none,                    
    numbersep=5pt,                  
    showspaces=false,                
    showstringspaces=false,
    showtabs=false,                  
    tabsize=2,
    frame=single,
    rulecolor=\color{gray}
}

\lstset{style=mystyle}

% Title Information
\title{\textbf{TensorFlow RTX 50 Series --- Conda + Docker Setup}}
\author{}
\date{}

\begin{document}

\maketitle

\tableofcontents
\newpage

\section{Project Overview}

\noindent This repository contains a complete configuration for running a \textbf{deep learning neural network} implemented in TensorFlow and trained on \textbf{NVIDIA GPUs} using \textbf{CUDA}. The application can run both on \textbf{GPU} and on \textbf{CPU}: at runtime, the logic in \texttt{main.py} automatically detects the available devices and decides whether to execute on a GPU (if present) or fall back to the CPU, printing a clear message in both cases.

\vspace{0.7em}

\noindent The project is designed to run efficiently on modern GPUs, including the \textbf{RTX 50 Series (5060, 5070, 5080, 5090)}, which are not officially supported by the standard TensorFlow GPU wheels at the time of development (2025).

\vspace{0.7em}

\noindent During the development of this application, TensorFlow did not provide precompiled kernels for \texttt{compute capability 12.0}. For this reason, the project relies on a \textbf{custom TensorFlow wheel} compiled specifically with RTX 50 support.

\vspace{0.7em}

\noindent The repository supports \textbf{two ways} of running the project:

\begin{itemize}
    \item[\textbf{1.}] \textbf{Conda Environment (Local Development)} --- runs directly on your host system and requires NVIDIA drivers, CUDA, and cuDNN installed on the host if you want GPU acceleration.
    \item[\textbf{2.}] \textbf{Docker Container (Reproducible Environment)} --- runs inside a container and can be executed either:
    \begin{itemize}
        \item on an \textbf{NVIDIA GPU} (using the NVIDIA Container Toolkit, without manually installing CUDA/cuDNN on the host), or
        \item on \textbf{CPU only}, without GPU acceleration and without any NVIDIA drivers.
    \end{itemize}
\end{itemize}

\noindent In the Conda method, the script \texttt{install\_tf.py} is used to detect the
available GPU and install the appropriate TensorFlow build (custom RTX 50
wheel or standard upstream wheel).
In the Docker method, TensorFlow is installed directly inside the image
at build time via the Dockerfile, and \texttt{install\_tf.py} is not used.

\subsection{CPU/GPU Device Selection in \texttt{main.py}}

\noindent The \texttt{main.py} file contains the logic that chooses the execution device at runtime. It uses TensorFlow to:

\begin{itemize}
    \item list the available physical GPU devices,
    \item select \texttt{/GPU:0} if at least one GPU is present,
    \item otherwise fall back to \texttt{/CPU:0},
    \item run a small matrix multiplication demo to verify that everything works correctly,
    \item print informative messages so the user knows whether the application is running on GPU or CPU.
\end{itemize}


\noindent This means that, once TensorFlow is installed correctly (via Conda or Docker), the application will automatically use the best available device without requiring manual configuration from the user.

\hrulefill

\section{Method 1 --- Conda Environment (Local Development)}

This method is intended for local development on your machine, using your host operating system. Here, GPU support depends on the \textbf{NVIDIA driver}, \textbf{CUDA}, and \textbf{cuDNN} installed on the host. If you only want to run on CPU, you can skip the CUDA/cuDNN installation and use a CPU-only TensorFlow build.

\subsection{System Requirements: NVIDIA Drivers, CUDA and cuDNN (for GPU)}

\subsubsection*{NVIDIA Drivers}

To use the GPU on Linux or WSL2, you must have \textbf{NVIDIA Drivers 560+} installed (or newer, such as the 580 series).

\begin{lstlisting}[language=bash]
nvidia-smi
\end{lstlisting}

\subsubsection*{System-Level CUDA and cuDNN}

\textit{Note: When running with the Conda method and using GPU acceleration, TensorFlow requires CUDA and cuDNN to be available at the \textbf{system level}, not only inside the virtual environment or Conda environment.}

\vspace{0.5em}

\noindent \textbf{Install CUDA 12.8 and cuDNN 9 (Ubuntu example):}
\begin{lstlisting}[language=bash]
sudo apt update
sudo apt install wget gnupg

# Add NVIDIA Repository
wget https://developer.download.nvidia.com/compute/cuda/repos/ubuntu2404/x86_64/cuda-ubuntu2404.pin
sudo mv cuda-ubuntu2404.pin /etc/apt/preferences.d/cuda-repository-pin-600

sudo apt-key adv --fetch-keys \
  https://developer.download.nvidia.com/compute/cuda/repos/ubuntu2404/x86_64/3bf863cc.pub

sudo add-apt-repository \
  "deb https://developer.download.nvidia.com/compute/cuda/repos/ubuntu2404/x86_64/ /"

# Install Toolkit and cuDNN
sudo apt update
sudo apt install -y cuda-toolkit-12-8
sudo apt install -y cudnn9-cuda-12
\end{lstlisting}

\noindent \textbf{Update PATH and Libraries:}
\begin{lstlisting}[language=bash]
echo 'export PATH=/usr/local/cuda-12.8/bin:$PATH' >> ~/.bashrc
echo 'export LD_LIBRARY_PATH=/usr/local/cuda-12.8/lib64:$LD_LIBRARY_PATH' >> ~/.bashrc
source ~/.bashrc
\end{lstlisting}

\subsection{Install Conda}

If you do not have Miniconda installed:

\begin{lstlisting}[language=bash]
wget https://repo.anaconda.com/miniconda/Miniconda3-latest-Linux-x86_64.sh -O miniconda.sh
bash miniconda.sh
source ~/.bashrc
\end{lstlisting}

\subsection{Create the Environment}

\begin{lstlisting}[language=bash]
conda env create -f environment.yml
conda activate tf-project
\end{lstlisting}

\subsection{TensorFlow Installation via \texttt{install\_tf.py}  (Local Only)}


For   \texttt{local Conda/venv environments}, this repository includes the script
\texttt{install\_tf.py}, which centralizes the TensorFlow installation logic. It
performs the following steps:

\begin{itemize}
    \item Detects the installed GPU model via \texttt{nvidia-smi}.
    \item If an \textbf{RTX 50 Series GPU} is detected $\rightarrow$ installs the custom TensorFlow wheel from \href{https://github.com/mypapit/tensorflowRTX50/releases}{mypapit/tensorflowRTX50}.
    \item If a different (officially supported) NVIDIA GPU is detected $\rightarrow$ installs the official TensorFlow GPU build.
    \item If no compatible GPU is found $\rightarrow$ installs a CPU-only TensorFlow build.
\end{itemize}

\noindent \textbf{Run the installer:}
\begin{lstlisting}[language=bash]
python install_tf.py
\end{lstlisting}

\textbf{Note:} \texttt{install\_tf.py} is intended \textbf{only for local environments}
(Conda or virtualenv). Docker images install TensorFlow directly via the
\texttt{Dockerfile}.

\subsection{Start the Project}

After TensorFlow is correctly installed, you can run the main application:

\begin{lstlisting}[language=bash]
python main.py
\end{lstlisting}

\hrulefill
\section{Method 2 --- Docker (Recommended)}

\emph{Recommended for grading, reproducible experiments, or users who do not
want to manually install CUDA and cuDNN on the host system.}

Docker allows you to isolate the environment and makes it easy to
reproduce the same configuration on different machines.

In the Docker setup:
\begin{itemize}
  \item TensorFlow is installed \textbf{at build time} inside the image.
  \item The script \texttt{install\_tf.py} is \textbf{not used inside the container}.
  \item A build argument (\texttt{TF\_VARIANT}) is used to choose between two image variants:
    \begin{itemize}
      \item a \textbf{custom RTX 50 image} (using the custom wheel),
      \item a \textbf{generic / official TensorFlow image} (for older GPUs or CPU-only).
    \end{itemize}
\end{itemize}

\subsection{Prerequisites}

\subsubsection*{Docker Installation}

You must have Docker (or Docker Desktop + WSL2 integration) installed on
your system. Please refer to the official Docker documentation for your
platform.

\subsubsection*{NVIDIA Container Toolkit (for GPU mode)}

If you want to run the container on an NVIDIA GPU, install the
\textbf{NVIDIA Container Toolkit} on the host:

\begin{lstlisting}[language=bash]
sudo apt install nvidia-container-toolkit
sudo nvidia-ctk runtime configure --runtime=docker
sudo systemctl restart docker
\end{lstlisting}

On Windows with Docker Desktop + WSL2, make sure:
\begin{itemize}
  \item GPU support is enabled in Docker Desktop,
  \item WSL2 integration is enabled for your Ubuntu distribution.
\end{itemize}

Once configured, you can run containers with GPU access using
\verb|--gpus all|.

\subsection*{Dockerfile Overview}

The provided \texttt{Dockerfile}:
\begin{itemize}
  \item Uses \verb|nvidia/cuda:12.8.0-cudnn-runtime-ubuntu24.04| as base image
        (CUDA + cuDNN already included).
  \item Creates a Python virtual environment at \verb|/opt/venv|.
  \item Installs base project dependencies from \texttt{requirements.txt}.
  \item Uses a build argument \texttt{TF\_VARIANT} to decide which TensorFlow
        variant to install:
        \begin{itemize}
          \item \verb|TF_VARIANT=rtx50| $\rightarrow$ installs the custom RTX 50
                wheel and extra dependencies from \texttt{requirements-rtx50.txt}.
          \item \verb|TF_VARIANT=official| $\rightarrow$ installs the official
                TensorFlow release from PyPI.
        \end{itemize}
  \item Sets \verb|CMD ["python", "main.py"]| so the demo runs automatically
        when the container starts.
\end{itemize}

\subsection{Build the Container Images}

From the root of the repository (where the \texttt{Dockerfile} is located), you
can build two images.

\subsubsection{1. Custom RTX 50 Image (for RTX 5060 / 5070 / 5080 / 5090)}


\begin{lstlisting}[language=bash]
docker build -t tf-app-rtx50 --build-arg TF_VARIANT=rtx50 .
\end{lstlisting}


This image:
\begin{itemize}
  \item installs the custom TensorFlow RTX 50 wheel,
  \item installs additional dependencies from \texttt{requirements-rtx50.txt},
  \item is intended for RTX 50 series GPUs when run with \verb|--gpus all|.
\end{itemize}

\subsubsection{2. Official TensorFlow Image (for older GPUs or CPU)}

\begin{lstlisting}[language=bash]
docker build -t tf-app-official --build-arg TF_VARIANT=official .
\end{lstlisting}


This image:
\begin{itemize}
  \item installs the official TensorFlow release from PyPI,
  \item can run on:
    \begin{itemize}
      \item older RTX GPUs (when run with \verb|--gpus all|),
      \item CPU-only environments (when run without GPU flags).
    \end{itemize}
\end{itemize}

\subsection{Running the Container on GPU (NVIDIA)}

To run the project using an NVIDIA GPU:


\begin{lstlisting}[language=bash]
docker run --gpus all -it tf-app-rtx50
\end{lstlisting}


or, for the official TensorFlow image:

\begin{lstlisting}[language=bash]
docker run --gpus all -it tf-app-official
\end{lstlisting}


In both cases:
\begin{itemize}
  \item The container has access to your host GPU(s) via \verb|--gpus all|.
  \item TensorFlow will attempt to use the GPU; if everything is configured
        correctly, \texttt{main.py} should detect at least one GPU and run on
        \verb|/GPU:0|.
\end{itemize}

\subsection{Running the Container on CPU Only}

If you do not have a compatible NVIDIA GPU, or if you want to test
CPU-only execution, you can start the official image without GPU flags:


\begin{lstlisting}[language=bash]
docker run -it tf-app-official
\end{lstlisting}

Behavior:
\begin{itemize}
  \item No GPU is exposed to the container.
  \item TensorFlow will run purely on CPU.
  \item \texttt{main.py} will print a message indicating that no GPU was found
        and it is running on \verb|/CPU:0|.
\end{itemize}



\hrulefill

\section{Included Files}

\begin{description}[style=multiline, leftmargin=4cm, font=\bfseries]
    \item[main.py] Main entry point for the deep learning application (model creation, training, evaluation, etc.). It also selects the execution device (GPU or CPU) at runtime and runs a small demo computation.
    \item[install\_tf.py] Hardware-aware installer script. Detects GPU type and installs the appropriate TensorFlow wheel (custom RTX 50 or official build; GPU or CPU-only).
    \item[environment.yml] Conda environment definition used for local development (Method 1).
    \item[Dockerfile] Build instructions for the Docker image used in Method 2 (GPU and CPU modes).
    \item[requirements.txt] Python dependencies used in the Docker image and for pip-based installations.
    \item[README.md] High-level project documentation and usage instructions.
\end{description}

\section{Final Notes}

\begin{itemize}
    \item At the time of development, the \textbf{RTX 50 Series} was not officially supported by standard TensorFlow GPU wheels, so a \textbf{custom build} is required.
    \item  TensorFlow installation, adapting to:
    \begin{itemize}
        \item RTX 50 Series GPUs (custom wheel),
        \item Other NVIDIA GPUs (official GPU wheel),
        \item CPU-only environments (no GPU detected).
    \end{itemize}
    \item The \textbf{Conda method} is suitable for local development but requires NVIDIA drivers, CUDA, and cuDNN on the host if you want GPU acceleration.
    \item The \textbf{Docker method} is recommended for reproducibility and for users who want to avoid manual system-level CUDA/cuDNN setup. It can be run on:
    \begin{itemize}
        \item GPU (via NVIDIA Container Toolkit and \verb|--gpus all|), or
        \item CPU only (no special host configuration required).
    \end{itemize}
    \item For questions, suggestions, or issues, please open an Issue in the repository.
\end{itemize}

\end{document}
